%%%%%%%%%%%%%%%%%%%%%%%%%%%%%%%%%%%%%%%%%%%%%%%%%%%%%%%%%%%%%%%%%%%%%%
% Lớp 4 | Chủ đề 1. Ôn tập và bổ sung | Ôn tập các số đến 100 000
% Dạng 4D1?1-2: Lập số từ các chữ số cho trước (lớn nhất, nhỏ nhất, đếm số lập được)
%%%%%%%%%%%%%%%%%%%%%%%%%%%%%%%%%%%%%%%%%%%%%%%%%%%%%%%%%%%%%%%%%%%%%%


% nguồn: 2026-11-18_lop4_chu-de-1-tu-luan (nhập 18/11/2026)
\begin{bt}%[4D1N1-2]%[ma:43dfbf2b]
% nguon: Toán tư duy lớp 4 – Bài 1 SGK Toán 4 KNTT
	Từ ba chữ số $3$, $7$, $0$, hãy viết tất cả các số có ba chữ số khác nhau.
	\loigiai{
		Chữ số hàng trăm phải khác $0$ nên chỉ có thể là $3$ hoặc $7$.\\
		Hàng trăm là $3$: được $307$; $370$.\\
		Hàng trăm là $7$: được $703$; $730$.\\
		Vậy có $4$ số: $307$, $370$, $703$, $730$.
	}
\end{bt}

% nguồn: 2026-11-18_lop4_chu-de-1-tu-luan (nhập 18/11/2026)
\begin{bt}%[4D1H1-2]%[ma:86372da0]
% nguon: Toán tư duy lớp 4 – Bài 1 SGK Toán 4 KNTT
	Từ năm chữ số $2$, $5$, $8$, $1$, $6$, hãy viết số lớn nhất và số nhỏ nhất có năm chữ số khác nhau. Tính hiệu của hai số đó.
	\loigiai{
		Số lớn nhất: xếp các chữ số từ lớn đến bé: $86\,521$.\\
		Số nhỏ nhất: xếp các chữ số từ bé đến lớn: $12\,568$.\\
		Hiệu: $86\,521 - 12\,568 = 73\,953$.
	}
\end{bt}

% nguồn: 2026-11-18_lop4_chu-de-1-tu-luan (nhập 18/11/2026)
\begin{bt}%[4D1V1-2]%[ma:f14a2b99]
% nguon: Toán tư duy lớp 4 – Bài 1 SGK Toán 4 KNTT
	Từ bốn chữ số $1$, $2$, $3$, $4$ có thể lập được bao nhiêu số có bốn chữ số khác nhau? Trong đó có bao nhiêu số chẵn?
	\loigiai{
		Hàng nghìn có $4$ cách chọn, hàng trăm còn $3$ cách, hàng chục còn $2$ cách, hàng đơn vị còn $1$ cách.\\
		Số các số lập được: $4 \times 3 \times 2 \times 1 = 24$ (số).\\
		Số chẵn có chữ số hàng đơn vị là $2$ hoặc $4$.\\
		Nếu hàng đơn vị là $2$: ba chữ số còn lại $1, 3, 4$ xếp được $3 \times 2 \times 1 = 6$ số. Tương tự hàng đơn vị là $4$: $6$ số.\\
		Vậy có $24$ số, trong đó có $6 + 6 = 12$ số chẵn.
	}
\end{bt}

% nguồn: 2026-11-18_lop4_chu-de-1-tu-luan (nhập 18/11/2026)
\begin{bt}%[4D1C1-2]%[ma:085bb172]
% nguon: Toán tư duy lớp 4 – Bài 1 SGK Toán 4 KNTT
	Có bao nhiêu số có ba chữ số mà tổng các chữ số bằng $5$?
	\loigiai{
		Chữ số hàng trăm khác $0$ và không lớn hơn $5$. Ta xét từng trường hợp hàng trăm:\\
		Hàng trăm $5$: chục + đơn vị $= 0$: số $500$ (1 số).\\
		Hàng trăm $4$: chục + đơn vị $= 1$: $410$, $401$ (2 số).\\
		Hàng trăm $3$: chục + đơn vị $= 2$: $320$, $311$, $302$ (3 số).\\
		Hàng trăm $2$: chục + đơn vị $= 3$: $230$, $221$, $212$, $203$ (4 số).\\
		Hàng trăm $1$: chục + đơn vị $= 4$: $140$, $131$, $122$, $113$, $104$ (5 số).\\
		Tổng cộng: $1 + 2 + 3 + 4 + 5 = 15$ (số).
	}
\end{bt}
